
\begin{remarques}
\label{XIII.1.5}
a) Sea~$F$ un haz de grupos sobre~$X$ tal que~$(F,f)$ sea
cohomológicamente propio relativamente a~$S$ en dimensión
$\leqslant -1$ (\resp $\leqslant 0$). Si se considera~$F$ como haz de
conjuntos, $(F,f)$ es cohomológicamente propio relativamente a~$S$ en
dimensión $\leqslant -1$ (\resp $\leqslant 0$), pero el recíproco es
falso.

Sea, por ejemplo, $Y$ el espectro de un anillo de valoración discreta
estrictamente local, con punto cerrado~$t$ y punto genérico~$s$;
sea $f\colon X\to Y$ un esquema no vacío sobre~$Y$ cuya fibra cerrada
sea vacía; sea~$F$ un haz constante no trivial de grupos sobre~$X$; y
sea~$P$ un torsor bajo~$F$ tal que
$\H^0(X_s,\leftexp{P\mkern-4mu}{F}|X_s)=1$. Entonces
$(\leftexp{P\mkern-4mu}{F},f)$ es cohomológicamente propio relativamente
a~$Y$ en dimensión $\leqslant -1$ cuando se considera
$\leftexp{P\mkern-4mu}{F}$ como haz de conjuntos. Si se considera
$\leftexp{P\mkern-4mu}{F}$ como haz de grupos, se tiene un isomorfismo
${\vphantom{F}}^{P^\circ\mkern-7mu}(\leftexp{P\mkern-4mu}{F})\simeq F$;
como el morfismo canónico
$$
\H^0(X,F)\to\H^0(X_t,F|X_t)=1
$$
no es inyectivo, esto demuestra que
$(\leftexp{P\mkern-4mu}{F},f)$ no es cohomológicamente propio
relativamente a~$Y$ en dimensión $\leqslant -1$.

b) Supongamos que~$f$ sea coherente (\ie cuasicompacto y
cuasiseparado). Sea~$F$ un stack sobre~$X$. Para todo punto geométrico
$\overline{y}$ de~$Y'$, se denota por~$\overline{Y}$ (\resp $\overline{Y'}$)
la localización estricta de~$Y$ (\resp $Y'$) en~$\overline{y}$, y se pone
$\overline{X}=X\times_Y\overline{Y}$,
$\overline{X'}=X'\times_{Y'}\overline{Y'}$, etc. Para que~$(F,f)$ sea
cohomológicamente propio relativamente a~$S$ en dimensión
$\leqslant -1$ (\resp $\leqslant 0$, \resp $\leqslant 1$), es necesario
y suficiente que, para todo $S$-esquema~$S'$ y todo punto geométrico
$\overline{y}$ de~$Y'$, el funtor canónico
$$
\overline{F}(\overline{X})\to\overline{F'}(\overline{X'})
$$
sea fiel (\resp plenamente fiel, \resp una equivalencia).

En
\marginpar{353}
efecto, si~$S'$ es un $S$-esquema, para que el funtor
$$
g^*f_*F\to f'_*F'
$$
sea fiel (\resp plenamente fiel, \resp una equivalencia), es necesario
y suficiente que lo sea el funtor inducido sobre las fibras en los
distintos puntos geométricos $\overline{y}'$ de~$\overline{Y'}$
\cite[III 2.1.5.9]{XIII.2}. La afirmación resulta, pues, del cálculo de
las fibras geométricas de la imagen directa de un stack por un morfismo
coherente \cite[VII 2.1.5]{XIII.2}.

c) Sea~$F$ un stack sobre~$X$. El hecho de que~$(F,f)$ sea
cohomológicamente propio relativamente a~$S$ en dimensión
$\leqslant -1$ (\resp $\leqslant 0$, \resp $\leqslant 1$) es local
sobre~$Y$ para la topología étale.

Sea~$S'$ un $S$-esquema, y sea~$F'$ la imagen inversa de~$F$ sobre~$X'$
(\cf \eqref{eq:XIII.1.0.1}). Si~$(F,f)$ es cohomológicamente propio
relativamente a~$S$ en dimensión $\leqslant 1$, también lo es~$(F',f')$.
Pero, si~$(F,f)$ es cohomológicamente propio relativamente a~$S$ en
dimensión $\leqslant -1$ (\resp $\leqslant 0$), no necesariamente sucede
lo mismo con~$(F',f')$.

Sea, por ejemplo, $S'$ un anillo de valoración discreta,
$f'\colon E_{S'}\to S'$ el espacio afín sobre~$S'$, $x$ un punto cerrado
de~$E_{S'}$ situado sobre el punto genérico de~$S'$, y~$F'$ el haz de
conjuntos sobre~$E_{S'}$ cuya restricción a~$E_{S'}-\{x\}$ es el haz
constante con un elemento y cuya fibra en un punto geométrico situado
sobre~$x$ tiene dos elementos. Entonces~$(F',f')$ no es
cohomológicamente propio relativamente a~$S'$ en dimensión
$\leqslant -1$. Sean $S=S'[Z]$, $f\colon E_S\to S$ el espacio afín
sobre~$S$ y~$T$ una parte cerrada de~$X=E_S$ que no interseca el cerrado
$Z=0$ y tal que~$f(T)$ contiene el punto genérico de~$S$.
Sean~$G$ la imagen inversa de~$F'$ sobre~$X$ y~$F$ el haz sobre~$X$
obtenido prolongando~$G|X-T$ por el vacío. Entonces~$(F,f)$ es
cohomológicamente propio relativamente a~$S$ en dimensión
$\leqslant -1$, pero deja de serlo tras el cambio de base $S'\to S$
definido por~$Z=0$.

d)
\marginpar{354}
Sea~$F$ un stack sobre~$X$ tal que~$(F,f)$ sea cohomológicamente propio
relativamente a~$Y$ en dimensión $\leqslant -1$ (\resp $\leqslant 0$,
\resp $\leqslant 1$). Entonces, para todo punto geométrico
$\overline{y}$ de~$Y$, el funtor canónico
$$
(f_*F)_{\overline{y}}\to F(X_{\overline{y}})
$$
es fiel (\resp plenamente fiel, \resp una equivalencia de categorías).
\end{remarques}
